%=============================================================
%  Chapter 4 --- State of the Art & Technology Choices
%=============================================================
\chapter{State of the Art and Technology Choices}

\section*{Introduction}
Building an intelligent e-commerce platform requires a rigorous selection of
technologies at every layer of the system.
This chapter presents a comparative study of existing solutions for each component
of \textbf{LUMINA}, justifies the choices made, and describes the hardware and
software development environment.

%-------------------------------------------------------------
\section{Frontend Framework}
%-------------------------------------------------------------

\subsection{Comparison}

The frontend is the first point of contact with the user.
The evaluation criteria are: rendering strategy (SSR/SSG), ecosystem maturity,
deployment simplicity, and TypeScript support.

\begin{table}[H]
\centering
\caption{Frontend framework comparison}
\begin{tabularx}{\textwidth}{|l|X|X|X|X|}
\hline
\rowcolor{gray!20}
\textbf{Criterion} & \textbf{Next.js 14} & \textbf{Nuxt 3} & \textbf{Remix} & \textbf{Create React App} \\
\hline
Server rendering (SSR/SSG) & \cellcolor{green!15}Native (App Router) & Native & Native & Limited \\
\hline
Routing & \cellcolor{green!15}File-based + layouts & File-based & File-based & Manual (React Router) \\
\hline
Performance & \cellcolor{green!15}Very high & High & High & Medium \\
\hline
Ecosystem & \cellcolor{green!15}Very large (Vercel) & Large (Vue) & Medium & Large \\
\hline
TypeScript & \cellcolor{green!15}Native & Native & Native & Optional \\
\hline
Image / Font optimization & \cellcolor{green!15}Built-in & Partial & Manual & Manual \\
\hline
Maturity & \cellcolor{green!15}Very mature & Mature & Recent & Deprecated \\
\hline
\end{tabularx}
\end{table}

\subsection{Choice: Next.js 14}

\textbf{Next.js 14} with the App Router was chosen for the following reasons:
\begin{itemize}
  \item Hybrid rendering (SSR for product page SEO, CSR for the cart and chatbot)
        is manageable at the component level through \textit{Server Components}
        and \textit{Client Components}.
  \item Automatic image optimization (\texttt{next/image}) is critical for a
        cosmetics store where visual quality is paramount.
  \item Co-location of routes, layouts, and loading states (\textit{loading.tsx},
        \textit{error.tsx}) reduces architectural complexity.
  \item TypeScript is supported natively, consistent with the rest of the stack.
\end{itemize}

%-------------------------------------------------------------
\section{Backend Framework}
%-------------------------------------------------------------

\subsection{Comparison}

\begin{table}[H]
\centering
\caption{Node.js backend framework comparison}
\begin{tabularx}{\textwidth}{|l|X|X|X|X|}
\hline
\rowcolor{gray!20}
\textbf{Criterion} & \textbf{Express 5 +\newline routing-controllers} & \textbf{NestJS} & \textbf{Fastify} & \textbf{Hono} \\
\hline
Architecture & \cellcolor{green!15}Flexible / Clean & Opinionated (MVC) & Flexible & Minimal \\
\hline
Performance & High & Medium & \cellcolor{green!15}Very high & Very high \\
\hline
Dependency injection & \cellcolor{green!15}TypeDI & Native (Nest) & Plugins & None \\
\hline
Decorators (TS) & \cellcolor{green!15}Yes (routing-controllers) & Yes & No & No \\
\hline
Ecosystem maturity & \cellcolor{green!15}Very mature & Mature & Mature & Recent \\
\hline
Learning curve & Low & High & Low & Very low \\
\hline
\end{tabularx}
\end{table}

\subsection{Choice: Express 5 + routing-controllers + TypeDI}

\textbf{Express 5} paired with \textbf{routing-controllers} and \textbf{TypeDI}
was preferred over NestJS despite surface-level similarities:
\begin{itemize}
  \item NestJS imposes a module structure that, for a project of this size,
        generates unnecessary verbosity and slows iteration.
  \item The Express + routing-controllers + TypeDI combination provides the same
        decorators (\texttt{@Controller}, \texttt{@Get}, \texttt{@Body},
        \texttt{@Inject}) without the NestJS framework overhead.
  \item Express 5 brings native async error handling, removing the need for
        \texttt{try/catch} wrappers in every route handler.
  \item The Clean Architecture applied to this project requires strict layer
        separation; an overly prescriptive framework like NestJS tends to
        short-circuit that separation.
\end{itemize}

%-------------------------------------------------------------
\section{ORM and Database}
%-------------------------------------------------------------

\subsection{Relational Database}

\begin{table}[H]
\centering
\caption{Database comparison}
\begin{tabularx}{\textwidth}{|l|X|X|X|}
\hline
\rowcolor{gray!20}
\textbf{Criterion} & \textbf{PostgreSQL} & \textbf{MySQL 8} & \textbf{MongoDB} \\
\hline
Model & \cellcolor{green!15}Relational + JSONB & Relational & Document (NoSQL) \\
\hline
ACID transactions & \cellcolor{green!15}Full & Full & Partial (since 4.0) \\
\hline
Complex queries & \cellcolor{green!15}Very strong (CTE, window fn.) & Good & Limited \\
\hline
Native JSON support & \cellcolor{green!15}JSONB + indexing & JSON (limited) & Native \\
\hline
Open source & \cellcolor{green!15}Yes & Yes (Oracle) & Yes (SSPL) \\
\hline
Production adoption & \cellcolor{green!15}Very wide & Wide & Wide \\
\hline
\end{tabularx}
\end{table}

\textbf{PostgreSQL} was chosen for its transactional robustness (critical for
checkout and stock management), its native \texttt{JSONB} type support
(used for the \texttt{skinType} and \texttt{ingredients} arrays on products),
and its full compatibility with Drizzle ORM.

\subsection{ORM}

\begin{table}[H]
\centering
\caption{TypeScript ORM comparison}
\begin{tabularx}{\textwidth}{|l|X|X|X|}
\hline
\rowcolor{gray!20}
\textbf{Criterion} & \textbf{Drizzle ORM} & \textbf{Prisma} & \textbf{TypeORM} \\
\hline
Type-safety & \cellcolor{green!15}Inferred (TS schema) & Generated (Prisma Client) & TS decorators \\
\hline
SQL control & \cellcolor{green!15}SQL-like API & Abstracted & Abstracted \\
\hline
Migrations & \cellcolor{green!15}drizzle-kit push/generate & prisma migrate & typeorm migration \\
\hline
Bundle size & \cellcolor{green!15}Very small & Medium (binary query engine) & Medium \\
\hline
Performance & \cellcolor{green!15}Very high & High & High \\
\hline
Maturity & Recent (stable) & \cellcolor{green!15}Mature & Mature \\
\hline
\end{tabularx}
\end{table}

\textbf{Drizzle ORM} was preferred over Prisma for two main reasons:
its SQL-like API enables complex filtered queries (joins, JSONB \texttt{@>}
operators for skin type filtering) without losing type safety;
and the absence of an external binary (Prisma's query engine) simplifies deployment.

%-------------------------------------------------------------
\section{Authentication}
%-------------------------------------------------------------

\begin{table}[H]
\centering
\caption{Authentication solution comparison}
\begin{tabularx}{\textwidth}{|l|X|X|X|X|}
\hline
\rowcolor{gray!20}
\textbf{Criterion} & \textbf{Better Auth} & \textbf{Auth.js (NextAuth)} & \textbf{Passport.js} & \textbf{Custom JWT} \\
\hline
Email OTP & \cellcolor{green!15}Native plugin & Partial & Manual & Manual \\
\hline
2FA (TOTP) & \cellcolor{green!15}Native plugin & No & Manual & Manual \\
\hline
Session + JWT & \cellcolor{green!15}Both & Sessions & JWT & JWT \\
\hline
Backend agnostic & \cellcolor{green!15}Yes & No (Next.js-centric) & Yes & Yes \\
\hline
JWT blacklist & \cellcolor{green!15}Via plugin / custom & No & Manual & Manual \\
\hline
TypeScript & \cellcolor{green!15}Native & Native & Partial & N/A \\
\hline
\end{tabularx}
\end{table}

\textbf{Better Auth} was chosen as it is the only framework offering native
\texttt{emailOTP} (signup verification) and \texttt{twoFactor} (TOTP/authenticator)
plugins without custom implementation.
Auth.js is tightly coupled to Next.js and does not integrate cleanly into a
separate Express backend --- a hard requirement of our architecture.
JWT blacklisting (token revocation on logout) is layered on top via Redis,
independently of Better Auth.

%-------------------------------------------------------------
\section{AI / RAG Stack}
%-------------------------------------------------------------

\subsection{Approach: RAG vs. Fine-tuning}

\begin{table}[H]
\centering
\caption{RAG vs. Fine-tuning for the product chatbot}
\begin{tabularx}{\textwidth}{|l|X|X|}
\hline
\rowcolor{gray!20}
\textbf{Criterion} & \textbf{RAG (Retrieval-Augmented Generation)} & \textbf{Fine-tuning} \\
\hline
Catalogue updates & \cellcolor{green!15}Immediate (re-indexing) & Full retraining \\
\hline
Infrastructure cost & \cellcolor{green!15}Low (local inference) & High (GPU training) \\
\hline
Response control & \cellcolor{green!15}Via indexed documents & Hard to control \\
\hline
Hallucinations & \cellcolor{green!15}Reduced (grounded context) & Frequent without sufficient data \\
\hline
Initial complexity & Medium & \cellcolor{green!15}Low (if cloud API) \\
\hline
Data required & \cellcolor{green!15}Catalogue documents & Thousands of examples \\
\hline
\end{tabularx}
\end{table}

The \textbf{RAG} approach is clearly suited to our context:
the product catalogue changes frequently, we do not have a sufficient training
dataset for quality fine-tuning, and grounding responses in real product sheets
is a business requirement.

\subsection{LLM Choice}

\begin{table}[H]
\centering
\caption{LLM comparison for the chatbot}
\begin{tabularx}{\textwidth}{|l|X|X|X|}
\hline
\rowcolor{gray!20}
\textbf{Criterion} & \textbf{Ollama (local)} & \textbf{OpenAI GPT-4} & \textbf{Mistral API} \\
\hline
Cost & \cellcolor{green!15}Free (self-hosted) & Paid (per token) & Paid (per token) \\
\hline
Data privacy & \cellcolor{green!15}Full (local) & Policy-dependent & Policy-dependent \\
\hline
Latency & Medium & \cellcolor{green!15}Low (global CDN) & Low \\
\hline
Response quality & Good (llama3, mistral) & \cellcolor{green!15}Excellent & Very good \\
\hline
Offline availability & \cellcolor{green!15}Yes & No & No \\
\hline
Integration & \cellcolor{green!15}Simple REST API & Official SDK & Official SDK \\
\hline
\end{tabularx}
\end{table}

\textbf{Ollama} (running locally, model \texttt{llama3}) was chosen for the
academic context of the project: no recurring cost, no customer data sent to
a third party, and a REST API identical to OpenAI's, facilitating a future
migration if needed.

\subsection{Embedding Model}

For semantic search, \textbf{sentence-transformers}
(\texttt{all-MiniLM-L6-v2}) was used: a compact model (22M parameters),
performant on short texts (product descriptions), and runnable locally without a GPU.

%-------------------------------------------------------------
\section{Cache Layer}
%-------------------------------------------------------------

\begin{table}[H]
\centering
\caption{Cache solution comparison}
\begin{tabularx}{\textwidth}{|l|X|X|X|}
\hline
\rowcolor{gray!20}
\textbf{Criterion} & \textbf{Redis} & \textbf{Memcached} & \textbf{In-process (Map)} \\
\hline
Data structures & \cellcolor{green!15}Rich (String, List, Set, ZSet) & Key-value only & Key-value only \\
\hline
Persistence & \cellcolor{green!15}Optional (RDB/AOF) & No & No \\
\hline
Native TTL & \cellcolor{green!15}Yes & Yes & Manual \\
\hline
Multi-instance & \cellcolor{green!15}Yes (cluster) & Yes & No \\
\hline
LUMINA use cases & \cellcolor{green!15}JWT blacklist, rate-limit, sessions & Limited & Limited \\
\hline
\end{tabularx}
\end{table}

\textbf{Redis} is used for three complementary use cases:
\begin{enumerate}
  \item \textbf{JWT blacklist} --- revoked token \texttt{jti} identifiers are stored
        with a TTL equal to the token's remaining lifetime.
  \item \textbf{Rate-limiting} --- sliding window per IP/email to protect against
        brute-force attacks on the login endpoint.
  \item \textbf{Better Auth sessions} --- optional server-side session storage.
\end{enumerate}

%-------------------------------------------------------------
\section{Development Environment}
%-------------------------------------------------------------

\subsection{Hardware}

\begin{table}[H]
\centering
\caption{Development machine specifications}
\begin{tabular}{|l|l|}
\hline
\rowcolor{gray!20}
\textbf{Component} & \textbf{Specification} \\
\hline
Operating System & Windows 11 Home \\
\hline
Processor & Intel Core i7 (12\textsuperscript{th} generation) \\
\hline
RAM & 16 GB DDR5 \\
\hline
Storage & NVMe SSD 512 GB \\
\hline
GPU & NVIDIA RTX (local Ollama inference) \\
\hline
Network & High-speed broadband (100 Mbps) \\
\hline
\end{tabular}
\end{table}

\subsection{Software}

\begin{table}[H]
\centering
\caption{Tools and versions used}
\begin{tabular}{|l|l|l|}
\hline
\rowcolor{gray!20}
\textbf{Tool} & \textbf{Version} & \textbf{Role} \\
\hline
Node.js & 20 LTS & Backend runtime \\
\hline
TypeScript & 5.x & Primary language (client + server) \\
\hline
Next.js & 14.x & Frontend framework \\
\hline
Express & 5.x & Backend HTTP framework \\
\hline
Drizzle ORM & 0.30.x & Database access \\
\hline
PostgreSQL & 16 & Primary database \\
\hline
Redis & 7.x & Cache / JWT blacklist \\
\hline
Python & 3.11 & RAG service \\
\hline
FastAPI & 0.111.x & RAG service API \\
\hline
Ollama & 0.1.x & Local LLM server \\
\hline
Docker & 25.x & Service containerization \\
\hline
VS Code & 1.90.x & Primary editor \\
\hline
Git / GitHub & 2.x & Version control \\
\hline
Postman & 11.x & REST API testing \\
\hline
PlantUML & 1.2024.x & UML diagram generation \\
\hline
Overleaf & -- & Report writing (LaTeX) \\
\hline
\end{tabular}
\end{table}

\subsection{Repository Structure}

The source code is organized as a monorepo with two root folders:

\begin{verbatim}
lumina/
+-- client/          # Next.js 14 - frontend
|   +-- app/         # App Router (pages, layouts)
|   +-- components/  # Reusable React components
|   +-- lib/         # API client, Zustand stores
+-- server/          # Express 5 - backend
|   +-- src/
|   |   +-- api/         # Controllers, DTOs, middlewares
|   |   +-- core/        # Use cases, entities, interfaces
|   |   +-- adapters/    # Repository implementations
|   |   +-- infrastructure/ # External services, DB, Redis
|   +-- drizzle/     # Schemas and migrations
+-- rag-service/     # Python FastAPI - RAG service
    +-- main.py
    +-- pipeline/
    +-- models/
\end{verbatim}

\section*{Conclusion}
This chapter presented and justified all technology choices for \textbf{LUMINA}.
The chosen stack --- Next.js 14, Express 5, Drizzle/PostgreSQL, Better Auth,
Redis, and Ollama --- forms a coherent ecosystem, fully typed in TypeScript,
and designed to meet the project's performance, security, and AI requirements.
The following chapters detail the concrete implementation of these technologies
across the three Scrum iterations.
